% =====================================================================
% Подробный разбор ex07: сервис локации и газовых алертов в шахте
%
% Сборка:
%   pdflatex main
%   pdflatex main   (повторно для оглавления)
%
% Требования:
%   - pdflatex (texlive-full / MiKTeX)
%   - Шрифты с кириллицей (любой современный TeX-дистрибутив)
% =====================================================================

\documentclass[11pt,a4paper]{article}

% ── Кириллица ────────────────────────────────────────────────────────
\usepackage[utf8]{inputenc}
\usepackage[T2A]{fontenc}
\usepackage[russian,english]{babel}

% ── Типографика и геометрия ──────────────────────────────────────────
\usepackage{geometry}
\geometry{margin=2.5cm}
\usepackage{parskip}
\usepackage{enumitem}
\setlist{noitemsep, topsep=4pt, partopsep=0pt}
\usepackage{amsmath}    % для \text{} в math mode

% ── Гиперссылки и оглавление ─────────────────────────────────────────
\usepackage{hyperref}
\hypersetup{
  colorlinks=true,
  linkcolor=blue!50!black,
  urlcolor=blue!50!black,
  citecolor=blue!50!black,
  pdftitle={miniFlink ex07: локация и газовые алерты},
  pdfauthor={miniFlink contributors}
}

% ── Код: listings для OCaml ──────────────────────────────────────────
% Настройка пакета listings для подсветки OCaml.
% Минимальный stylesheet, работает без внешних зависимостей.

\usepackage{listings}
\usepackage{xcolor}

% Цвета (мягкие, нейтральные)
\definecolor{kwcolor}{rgb}{0.20,0.20,0.65}   % ключевые слова — тёмно-синий
\definecolor{strcolor}{rgb}{0.30,0.55,0.25}  % строки — зелёный
\definecolor{cmtcolor}{rgb}{0.45,0.45,0.45}  % комментарии — серый
\definecolor{bgcolor}{rgb}{0.98,0.98,0.96}   % фон листинга

% Язык OCaml не встроен в listings — определяем словарь сами
\lstdefinelanguage{OCaml}{
  keywords={
    let, in, fun, function, match, with, when, type, and, or,
    if, then, else, begin, end, module, struct, sig, val,
    open, include, of, rec, mutable, ref, true, false,
    while, do, done, for, to, downto, exception, try, raise,
    assert, lazy, object, class, method, inherit, new, virtual,
    private, initializer
  },
  keywordstyle=\color{kwcolor}\bfseries,
  sensitive=true,
  comment=[s]{(*}{*)},
  morecomment=[s][\color{cmtcolor}\itshape]{(*}{*)},
  string=[b]",
  stringstyle=\color{strcolor},
}

\lstset{
  language=OCaml,
  basicstyle=\ttfamily\small,
  backgroundcolor=\color{bgcolor},
  frame=single,
  framerule=0.4pt,
  rulecolor=\color{gray!40},
  xleftmargin=8pt,
  xrightmargin=8pt,
  framexleftmargin=4pt,
  framexrightmargin=4pt,
  framextopmargin=2pt,
  framexbottommargin=2pt,
  showstringspaces=false,
  breaklines=true,
  breakatwhitespace=true,
  tabsize=2,
  captionpos=b,
  aboveskip=8pt,
  belowskip=8pt,
  numbers=none,
  literate=
    {→}{$\to$}1 {←}{$\leftarrow$}1
    {≥}{$\geq$}1 {≤}{$\leq$}1 {≠}{$\neq$}1
    {÷}{$\div$}1 {±}{$\pm$}1,
  extendedchars=true,
  inputencoding=utf8,
}

% Для ASCII-диаграмм — отдельный стиль без подсветки
\lstdefinestyle{diagram}{
  language=,
  basicstyle=\ttfamily\footnotesize,
  frame=none,
  backgroundcolor=\color{bgcolor},
  keywordstyle=,
}


% ── Графика (для ASCII-диаграмм будем использовать verbatim) ─────────
\usepackage{xcolor}
\definecolor{boxgray}{rgb}{0.95,0.95,0.95}
\usepackage{tcolorbox}
\tcbset{
  colback=boxgray, colframe=gray!50,
  boxrule=0.5pt, arc=2pt,
  left=8pt, right=8pt, top=4pt, bottom=4pt,
}

% =====================================================================
\title{\textbf{miniFlink ex07}\\[4pt]
  \large Сервис локации и газовых алертов для подземной шахты:\\
  пошаговый разбор архитектуры}
\author{}
\date{}

\begin{document}

\maketitle

\begin{abstract}
\noindent
ex07 --- самый большой пример в репозитории miniFlink. Он строит сервис
для подземной угольной шахты, который из двух независимых потоков
(пакеты от фонарей и пакеты от газовых датчиков) выдаёт три типа
выходов: алерты состояния шахтёра, локацию по маякам связи и
обогащённые координатами газовые алерты. Под капотом используется
почти весь арсенал библиотеки: скользящие окна с
\texttt{allowed\_lateness}, \texttt{group\_by}, FSM через
\texttt{process\_keyed}, объединение потоков с выравниванием
watermark, retract-семантика для обновлений.

Эта статья --- подробный разбор того, \textit{зачем} каждое решение
принято таким, как оно принято: какие альтернативы рассматривались,
почему отброшены, что бы случилось при наивной реализации. Цель ---
дать инженеру, выбирающему stream processing, конкретный материал для
оценки, а контрибьютору~--- понимание design-rationale.
\end{abstract}

\tableofcontents
\bigskip

% ─────────────────────────────────────────────────────────────────────
% ─────────────────────────────────────────────────────────────────────
% Раздел 1 — Задача
% ─────────────────────────────────────────────────────────────────────

\section{Задача}

Представьте себе подземную угольную шахту: километры тоннелей на
нескольких горизонтах, в среднем по сорок-сто шахтёров на смену, у
каждого фонарь, который раз в пятнадцать секунд шлёт радиопакет с
показаниями нескольких ближайших маяков связи, а раз в десять секунд
--- отдельный пакет с показаниями газовых сенсоров. Эти пакеты летят
в диспетчерскую через смесь UDP-радио и LTE-шлюзов, доходят с
задержками, иногда дублируются, иногда теряются.

В диспетчерской на стене висит экран с картой шахты и состоянием
смены: где находится каждый шахтёр, у кого тревога по газу, кто
прижал кнопку SOS, у кого вот-вот сядет батарейка фонаря. Этот экран
обслуживает наш сервис.

\subsection{Что должен делать сервис}

Если разложить требования \emph{диспетчера} на конкретные пайплайны,
получается три независимых подсистемы:

\begin{description}
\item[Контроль состояния шахтёра.] Из RSSI-пакета извлекаем не
показания маяков (это работа другой подсистемы), а \emph{косвенные
факты}: пакет вообще пришёл, в нём были показания, фонарь двигался,
напряжение батареи в норме. По каждому из этих фактов держим
независимый таймер и эмитим соответствующий алерт когда таймер
перевалил порог: ``не было пакетов больше двух минут'',
``пакеты идут, но маяков не слышит больше пяти'', ``не двигался дольше
порога'', ``батарея ниже 3{,}5 В с дребезгом по гистерезису''.
Плюс импульс SOS (нажатая кнопка) --- мгновенный без таймера.

\item[Локация по маякам.] Для каждого шахтёра в скользящем
60-секундном окне с шагом 5 секунд берём топ-2 маяков с самой сильной
медианной RSSI и из них приближённо вычисляем координаты шахтёра
через интерполяцию между двумя маяками пропорционально расстоянию.
Окно скользит каждые 5 секунд, поэтому диспетчерский экран обновляется
не реже чем раз в 5 с; при опаздавшем пакете окно
\emph{переоткрывается} и эмитит \texttt{Retract} старого значения с
последующим обновлённым \texttt{Data} --- так картинка остаётся
консистентной.

\item[Газовые алерты с обогащением координатами.] Газовый пакет
приходит \emph{отдельно} и сам по себе не несёт ID маяков или
позиции. Когда газ превышает порог --- алерт идёт немедленно
(безопасность важнее красоты), но без координат, если их ещё не было.
Как только локация для этого шахтёра становится известна (хотя бы по
одному маяку), алерт автоматически \emph{обновляется}: старая версия
без координат retract'ится, новая с координатами эмитится. То же самое
при смене уровня опасности (Warning $\to$ Critical) или
существенном изменении ppm. Когда газ возвращается в норму --- алерт
снимается через \texttt{Gas\_resolved}.
\end{description}

Эти три подсистемы независимы по бизнес-логике, но обмениваются
данными: газовый пайплайн \emph{потребляет} результат пайплайна
локации, чтобы обогатить свои алерты координатами.

\subsection{Что в этой задаче сложного}

Если бы это была учебная задача из курса по базам данных, её можно
было бы решить за полчаса: батчевый ETL раз в минуту, JOIN по
шахтёру, отрисовка в дашборд. Но в реальной шахте такое решение
\emph{бесполезно}, и причин этому несколько.

\paragraph{Время событий не равно времени поступления.} Пакет от
шахтёра M1 с временной меткой \texttt{t=88s} может прийти в
диспетчерскую в момент \texttt{t=210s}: UDP в шахте теряет пакеты,
LTE-шлюз делает повторные попытки, на агрегаторе очереди.
Если мы будем считать алерты по wall-clock времени поступления, то
получим ``M1 не двигался последние 2 минуты'' в момент когда M1 уже
давно дома --- его пакет о движении просто опаздывает по дороге.
Сервис должен работать по \emph{event-time} (времени самого
события), а не по wall-clock.

\paragraph{Идемпотентность.} Тот же UDP-механизм даёт дубликаты:
пакет может прийти дважды, потому что шлюз ретранслировал его при
сомнениях. Если мы будем считать тревогу по второму дублю
``независимо'', получим \emph{два} алерта SOS на одно нажатие кнопки.
Диспетчер не должен видеть фантомных тревог.

\paragraph{Опоздавшие пакеты ломают окна.} Окно
\texttt{[60s, 120s)} закрылось в 120с, и пайплайн эмитнул финальное
top-2 маяков. Но в 145с пришёл опоздавший пакет с
\texttt{ts=90s} --- он \emph{принадлежит} уже закрытому окну. Нужно
либо его выбросить (теряем данные), либо переоткрыть окно, пересчитать
агрегат и эмитнуть новое значение, помечая старое как недействительное.
Второе правильно для диспетчера; первое было бы дешевле, но даёт
непредсказуемый дашборд (на одной и той же временной метке могут быть
разные значения в разные моменты wall-clock).

\paragraph{Газовый алерт без координат лучше чем нет алерта.}
Газ начинает выделяться \emph{раньше} чем шахтёр успевает прийти под
маяк --- особенно в забое, где маяки ставят редко. Если ждать
координат прежде чем эмитить тревогу о 120 ppm CO, можно
потерять минуты на ожидании. Алерт идёт сразу. Но как только
координаты появляются --- они должны \emph{подтянуться} к уже
существующей тревоге, не порождая дубль и не давая диспетчеру повод
закрыть старый алерт и открыть новый вручную.

\paragraph{Объём.} Шахта на сорок шахтёров --- мелочь. Реальные
большие шахты в Кузбассе или Караганде имеют до четырёх тысяч
работников в смену. RSSI-пакетов от такой шахты получается порядка
270 ev/s, газовых --- ещё 400 ev/s. Это
немного по меркам стримов, но \emph{пайплайн должен держать нагрузку
с запасом}: пиковые моменты (несколько горящих участков, грозовая
помеха в радио, перезапуск шлюза) дают взрывы в десятки тысяч событий
в секунду на короткий период. Сервис не должен в эти моменты
терять алерты.

\subsection{Что есть готового в miniFlink}

Библиотека даёт нам почти всё что нужно для этого сервиса, без
необходимости писать stream-runtime самостоятельно:

\begin{itemize}
\item \texttt{Pipe.event\_time \textasciitilde lateness} ---
вставка watermark'ов в поток по event-time;
\item \texttt{Pipe.dedup} --- дедупликация по rule-функции;
\item \texttt{Pipe.window\_agg\_keyed} с
\texttt{allowed\_lateness} --- скользящие окна с retract'ами при
опоздавших пакетах;
\item \texttt{Agg.median}, \texttt{Agg.group\_by},
\texttt{Agg.top\_k\_by} --- набор комбинируемых агрегатов;
\item \texttt{Pipe.process\_keyed} --- per-key state с self-timers
(для FSM состояния шахтёра);
\item \texttt{Mf\_event.union} --- объединение потоков с
выравниванием watermark'а по минимуму (для трёх-поточного газового
пайплайна);
\item \texttt{Pipe.materialize} --- сворачивание retract-богатого
потока в финальную таблицу записей (для sink-а).
\end{itemize}

Чего \emph{нет} --- это бизнес-логики и доменных типов. Их пишем мы;
структура ниже как раз и есть рассказ о том, как мы их написали.

\subsection{Что увидим дальше}

Статья идёт от внешнего наружу. В разделе~\ref{sec:inputs}
разбираем устройство входов: почему два независимых потока вместо
одного, какие в них аномалии. В разделе~\ref{sec:pipelines} --- три
пайплайна по очереди: FSM состояния, скользящие окна с медианой и
интерполяцией позиции, retract-обогащение газовых алертов.
Раздел~\ref{sec:retract} собирает наблюдения о retract-семантике как
о сквозном механизме, проходящем через все три пайплайна.
Раздел~\ref{sec:performance} приводит измерения производительности
(на машине автора-эксперимента и на пользовательском i7-8700 с OCaml
5) и описывает направленные оптимизации, которые мы делали.
Раздел~\ref{sec:conclusions} --- что отложено, что не делалось
намеренно, и при каких условиях стоит реализовать недостающее.

% Раздел 2 — Входы: пайплайны источников
\section{Входы}\label{sec:inputs}
\textit{Этот раздел ещё не написан.}

% ─────────────────────────────────────────────────────────────────────
% Раздел 3 — Три пайплайна
% ─────────────────────────────────────────────────────────────────────

\section{Три пайплайна}\label{sec:pipelines}

Три пайплайна образуют ядро сервиса. У каждого свой набор операторов
и свои дизайнерские решения, но между ними виден общий стиль: все
три потребляют один и тот же \texttt{packet}-поток (или его
производные), все три используют event-time, у всех трёх результат
идёт в sink-обвязку, которая отображает или публикует данные.

Разбираем по одному. Самый простой по семантике, но самый
интересный архитектурно --- \texttt{connectivity\_alerts}.

\subsection{connectivity\_alerts: три слоя FSM}\label{ssec:connectivity}

Этот пайплайн ловит пять диагностических состояний шахтёра:

\begin{itemize}
\item \textbf{No\_packets} --- пакеты от фонаря перестали приходить
($> 2$ мин тишины);
\item \textbf{No\_readings} --- пакеты идут, но \texttt{readings}
пустой ($> 5$ мин не слышит ни одного маяка);
\item \textbf{No\_motion} --- датчик ускорения молчит ($> 2$ мин
без движения);
\item \textbf{Low\_voltage} --- напряжение батареи ниже 3{,}5 В,
держится дольше 2 минут (с гистерезисом возврата 3{,}7 В);
\item \textbf{Sos} --- кнопка SOS нажата (мгновенно, без таймера).
\end{itemize}

Любая из четырёх первых тревог снимается, когда условие перестаёт
выполняться (для напряжения --- с гистерезисом, об этом ниже).

\paragraph{Соблазн монолита.} Самое прямое решение --- один большой
\texttt{process\_keyed.on\_event} с per-key state, в котором
аккумулируем «время последнего пакета, время последнего движения,
текущее напряжение, флаг что был ли уже Low\_voltage». На каждом
пакете обновляем state, смотрим что изменилось, эмитим алерты.

Это работает, но плохо читается и плохо тестируется. На каждом
пакете мы одновременно делаем три разных вещи:
\textbf{запоминаем факты} (когда последний пакет),
\textbf{принимаем решения} (нужно ли алертить),
\textbf{регистрируем таймеры} (когда проверить тишину снова, если
сейчас пакетов нет). Все три смешаны в одной функции, и поправить
поведение для одного типа алерта означает читать весь файл.

\paragraph{Разделение на три слоя.} Поэтому
\texttt{Pipelines} раскладывает FSM на три модуля:

\begin{itemize}
\item \textbf{Last\_seen} --- хранилище фактов;
\item \textbf{Fsm} --- чистые функции «факты $\to$ алерт»;
\item \textbf{Self\_timer} --- регистрация таймеров.
\end{itemize}

\texttt{process\_keyed} склеивает их в общий пайплайн, но каждый
модуль изолирован и тестируем независимо.

\subsubsection{Слой 1: Last\_seen}

Просто запись того, что мы знаем о шахтёре в данный момент.

\begin{lstlisting}
type t = {
  mutable last_pkt_ts    : Time.t;   (* any packet *)
  mutable last_reading_ts: Time.t;   (* packet with non-empty readings *)
  mutable last_motion_ts : Time.t;   (* packet with moving=true *)
  mutable last_voltage   : float;
  mutable last_voltage_ts: Time.t;
  mutable low_v_since    : Time.t option;
  mutable in_low_v       : bool;     (* for debounce + hysteresis *)
}
\end{lstlisting}

Метод \texttt{record} обновляет нужные поля по приходящему пакету.
В нём \emph{нет} логики «что эмитить» --- только запоминание.

Этот слой решает проблему «когда таймер сработал в момент молчания,
у нас нет пакета чтобы взять из него факты» --- факты уже лежат в
\texttt{Last\_seen}.

\subsubsection{Слой 2: Fsm}

Чистые функции от \texttt{Last\_seen} к опциональному алерту:

\begin{lstlisting}
let no_packets seen ~now =
  if now - seen.last_pkt_ts > no_packets_threshold
  then Some (No_packets (seen.lamp, seen.last_pkt_ts))
  else None

let no_motion seen ~now = ...
let no_readings seen ~now = ...
let low_voltage seen ~now = ...

(* Aggregate: at what nearest moment in time should we
   re-check the state, assuming no new packet arrives? *)
let next_check seen = ...
\end{lstlisting}

В этом слое нет state и нет I/O --- любая функция здесь принимает
\texttt{Last\_seen.t} и возвращает чистое значение. Тесты пишутся
тривиально: «вот snapshot, вот ожидаемый алерт».

\paragraph{next\_check как сердце таймера.} Самая интересная функция
здесь --- \texttt{next\_check}. Она отвечает на вопрос: «если сейчас
\emph{не} придёт нового пакета, в какой ближайший момент состояние
может измениться?» Например, если последний пакет был в 100с и
порог тишины --- 2 минуты, то изменение наступит в 220с
(100 + 120). Если последнее движение было в 95с --- то в 215с. И
так далее. \texttt{next\_check} берёт минимум по всем «пробудкам»
и возвращает один timestamp.

Это нужно \texttt{Self\_timer}'у, чтобы знать, на какой момент
будить пайплайн.

\subsubsection{Слой 3: Self\_timer}

\texttt{process\_keyed} в miniFlink имеет встроенную поддержку
self-timers через \texttt{ctx.set\_timer}. \texttt{Self\_timer}
тонко оборачивает её, добавляя одну инвариантность: на каждом ключе
живой может быть \emph{максимум один} таймер. Когда приходит новый
пакет и \texttt{Fsm.next\_check} даёт новое время будильника, старый
таймер отменяется и регистрируется новый. Это уберегает от
накопления старых таймеров (особенно опасно для долгоживущих
шахтёров).

\begin{lstlisting}
let reschedule t ctx ~target =
  match t.armed with
  | Some old when old = target -> ()       (* same timer *)
  | _ ->
    (* Old timer will be ignored on fire -- we changed `armed`,
       so when it fires the handler sees it's stale. *)
    ctx.Pipe.set_timer target Pipe.Wallclock;
    t.armed <- Some target
\end{lstlisting}

\subsubsection{Склейка через process\_keyed}

\begin{lstlisting}
let connectivity_alerts source =
  source
  |> Pipe.event_time ~lateness:(seconds 1)
  |> Pipe.process_keyed (module ByLamp)
       ~init:make_state
       ~on_event:(fun ctx key st p ->
         (* Layer 1: remember facts *)
         (match Last_seen.record st.seen ~p with
          | `Sos_pressed -> ctx.emit (Sos (key, p.ts))
          | `Quiet -> ());
         (* Layer 2 + emission *)
         check_and_emit ctx key st ~now:p.ts;
         (* Layer 3: schedule the next check *)
         Self_timer.reschedule st.timer ctx
           ~target:(Fsm.next_check st.seen))
       ~on_timer:(fun ctx key st t _kind ->
         Self_timer.consumed st.timer;
         check_and_emit ctx key st ~now:t;
         let next = Fsm.next_check st.seen in
         if next > t && next < max_int then
           Self_timer.reschedule st.timer ctx ~target:next)
\end{lstlisting}

Двадцать строк кода, и каждая строка делает ровно одну вещь. Это
прямой эффект разделения: ни одна функция в этом пайплайне не
смешивает «помню что было» с «решаю что делать».

\paragraph{Гистерезис напряжения.} Один из неприметных, но важных
моментов --- \texttt{Low\_voltage} имеет \emph{два} порога:
\texttt{low\_voltage\_threshold = 3.5} (вход в подозрение) и
\texttt{voltage\_ok\_threshold = 3.7} (выход обратно). Это
\emph{гистерезис}: вошёл в Low\_voltage при $V < 3{,}5$, вышел только
при $V > 3{,}7$. Между ними --- \emph{мёртвая зона}, в которой
состояние не меняется.

Без гистерезиса колебания напряжения около 3{,}5 В (нормальное
поведение под нагрузкой фонарика) дали бы \emph{дребезг алертов} ---
Low\_voltage, Voltage\_ok, Low\_voltage, Voltage\_ok --- много раз в
минуту. Диспетчер увидел бы хаос. Один из тех мелких моментов,
которые отличают workable систему от непригодной.

Дополнительно есть \texttt{voltage\_debounce = 2 минуты}: алерт
\texttt{Low\_voltage} эмитится не сразу как напряжение упало ниже 3{,}5,
а только если оно \emph{держится} там 2 минуты подряд. Это уже не
гистерезис, это debounce --- защита от кратковременных просадок
напряжения (включение фонаря на полную мощность даёт мгновенный
провал).

\paragraph{Почему такой подход универсален.} Каждый раз, когда в
проде встаёт задача «детектировать состояние из временного ряда
событий», возникают те же три проблемы:
а) надо помнить факты;
б) надо принимать решения о них;
в) надо реагировать на молчание.
Three-layer split здесь не специфичен для шахты или для miniFlink ---
он будет применим к любой подобной задаче (мониторинг IoT-устройств,
heartbeats веб-сервисов, dropped sessions в чатах). В этом примере
он реализуется в идиомах OCaml + \texttt{process\_keyed}, в Java
+ Flink \texttt{KeyedProcessFunction} был бы тот же split в тех
же трёх классах.

\medskip

Дальше переходим к пайплайну локации --- архитектурно проще
(один длинный pipe из операторов), но с насыщенной алгоритмикой:
скользящие окна с retract'ами, медиана по биконам, top-2, и
линейная интерполяция координат.

\subsection{median\_rssi: скользящие окна и интерполяция координат}\label{ssec:median}

Задача пайплайна --- для каждого шахтёра выдавать «топ-2 маяков с
самой сильной медианной RSSI» в скользящем 60-секундном окне с
шагом 5 секунд, а из них приближённо вычислять координаты шахтёра в
двумерной плоскости горизонта.

Архитектурно это \emph{линейный} pipe из операторов; интересна не
структура, а каждый отдельный шаг и почему он выглядит так.

\subsubsection{Шаги пайплайна}

\paragraph{1. dedup по бизнес-ключу.}
\begin{lstlisting}
|> Pipe.dedup (module ByLamp)
     ~rule:(fun p -> string_of_int p.ts)
     ~cooldown:(seconds 120)
\end{lstlisting}

\texttt{Pipe.dedup} держит per-key cache недавно виденных rule-значений
и фильтрует повторы. Ключ дедупликации --- \texttt{(lamp, ts)},
выражаемый как \texttt{string\_of\_int p.ts} в пределах ключа
\texttt{p.lamp}. Если пакет от M1 с \texttt{ts=88s} пришёл дважды
(at-least-once Kafka, UDP retry), мы дедуплицируем по бизнес-ключу
\emph{пакета} (lamp + event-time), а не по байтовому содержимому
сетевого сообщения --- которое шлюз вполне мог обогатить полем
\texttt{received\_at} перед ретрансляцией.

\texttt{cooldown 120s} означает: дедуплицируем дубли в пределах двух
минут. Если тот же ключ придёт через 3 минуты --- это уже не дубль,
а что-то странное (битый clock на устройстве?), но пропускаем дальше,
не теряя данные.

\paragraph{2. flat\_map: разворот в reading'и.}
\begin{lstlisting}
|> Pipe.flat_map (fun p ->
     List.map (fun (b, rssi) ->
       { r_lamp = p.lamp; r_beacon = b;
         r_rssi = rssi; r_ts = p.ts })
       p.readings)
\end{lstlisting}

Каждый пакет имеет от одного до пяти reading'ов. \texttt{flat\_map}
разворачивает один пакет в несколько событий, каждое --- запись
вида ``шахтёр X слышал маяк Y с уровнем Z в момент T''. Дальше
работаем с плоским потоком, а не с пакетами.

Почему это правильно: для медианы по бикону нам нужно
\emph{сгруппировать} все RSSI данного маяка отдельно, что естественно
делается на плоских записях. Если бы оставили пакеты,
\texttt{group\_by} пришлось бы делать по \texttt{packet.readings}, и
агрегаты усложнились бы.

\paragraph{3. event\_time с lateness.}
\begin{lstlisting}
|> Pipe.event_time ~lateness:(seconds 1)
\end{lstlisting}

Вставка watermark'ов на основе event-time каждого reading'а. Параметр
\texttt{lateness:1s} означает: «watermark отстаёт от максимального
виденного event-time на 1 секунду». Это \emph{out-of-order tolerance}
между близкими во времени событиями: если у нас сейчас прошло событие
с \texttt{ts=100s}, watermark поднимется только до 99с, давая 1
секунду на ``последние опоздавшие'' до того как окно закроется.

\paragraph{Почему именно 1 секунда.} В нашем сценарии reading'и
сгенерированы из одного пакета с одинаковым \texttt{ts} --- они
синхронны. Out-of-order возникает только между пакетами от разных
шахтёров, доставленными через сеть. Эмпирически разброс доставки
для не-aномальных пакетов укладывается в 1с. Если поставить больше
(скажем, 60с) --- окна закрываются медленнее, дашборд лагает на
минуту. Меньше --- теряем близко-опоздавшие. Это типичный
out-of-order trade-off; 1с в нашем сценарии нормально, в шахтах с
очень плохим радио может потребоваться 5--10с.

Сама обработка \emph{сильно} опоздавших пакетов (более 60с) лежит на
\texttt{allowed\_lateness} следующего оператора --- они переоткрывают
окно через retract.

\paragraph{4. window\_agg\_keyed: главный оператор.}
\begin{lstlisting}
|> Pipe.window_agg_keyed
     ~by:(fun r -> r.r_lamp)
     ~allowed_lateness:(seconds 60)
     (Pipe.sliding (seconds 60) (seconds 5))
     Agg.(...)
\end{lstlisting}

Здесь несколько параметров:

\begin{itemize}
\item \texttt{by} --- ключ партиционирования (шахтёр);
\item \texttt{sliding (60s) (5s)} --- окно длиной 60с с шагом 5с,
то есть событие $t=30$с попадает в окна $[-30, 30)$, $[-25, 35)$,
$[-20, 40)$ \ldots, $[30, 90)$. Двенадцать окон на событие;
\item \texttt{allowed\_lateness 60s} --- если опоздавший пакет
попадает в \emph{уже закрытое} окно, но не более чем на 60с
прошлое --- окно \textbf{переоткрывается}, агрегат пересчитывается,
эмитится \texttt{Retract} старого результата плюс \texttt{Data}
нового. Пакеты опоздавшие сильнее --- идут в side output (мы их не
собираем; в проде шли бы в DLQ);
\item \texttt{Agg.t} --- комбинаторно-собираемый агрегат, см. ниже.
\end{itemize}

Это \emph{единственный} stateful оператор всего пайплайна. Состояние
--- хеш-таблица «окно $\to$ список значений (или сложенный
аккумулятор)», очищается автоматически при закрытии окна (или
оставляется на время \texttt{allowed\_lateness}).

\paragraph{5. Двухуровневая агрегация через Agg.}
\begin{lstlisting}
Agg.(
  group_by
    ~key:(fun r -> r.r_beacon)
    ~inner:(median (fun r -> r.r_rssi))
  |> map (fun per_beacon ->
       per_beacon
       |> List.filter_map (fun (b, med) ->
            match med with
            | Some m -> Some (b, m)
            | None -> None)
       |> Agg.run (top_k_by 2 ~by:snd)))
\end{lstlisting}

Внутри одного окна шахтёра М1 у нас может быть, например, 8
reading'ов от B1, 12 от B2, 5 от B3. \texttt{group\_by ~key:beacon}
группирует их по маяку. \texttt{\textasciitilde inner: median} ---
агрегат, применяемый к каждой группе: для B1 даёт medianX,
для B2 --- medianY, для B3 --- medianZ. На выходе списка пар
\texttt{[(B1, X); (B2, Y); (B3, Z)]}.

Дальше \texttt{map} превращает этот список в top-2 через
\texttt{top\_k\_by 2}. У нас два уровня агрегации --- сначала
медиана внутри маяка, потом ranking между маяками --- собранные
\emph{декларативно} из библиотечных кубиков, без рукописных циклов.

\subsubsection{Линейная интерполяция координат}\label{sssec:interpolate}

Имея top-2 маяков с известными координатами и медианной RSSI до
каждого, мы хотим приближённо найти координаты шахтёра в плоскости
горизонта.

\paragraph{От RSSI к расстоянию.} В \texttt{Large\_mine} RSSI
генерируется по формуле $\text{RSSI} = -40 - 20 \log_{10}(d)$.
Инверсия:
$$d = 10^{\frac{-\text{RSSI} - 40}{20}}$$

\begin{lstlisting}
let distance_from_rssi (rssi : float) : float =
  10. ** ((-. rssi -. 40.) /. 20.)
\end{lstlisting}

В реальной шахте этот калибровочный коэффициент зависит от частоты
радио, влажности воздуха, рельефа тоннеля, и в проде калибруется
эмпирически на каждом горизонте. У нас вшито в код --- демо.

\paragraph{Интерполяция между двумя маяками.} Имея A с координатами
\texttt{(xa, ya)} и расстоянием \texttt{d\_a}, B с координатами
\texttt{(xb, yb)} и расстоянием \texttt{d\_b}, точку на отрезке
$AB$ пропорционально расстояниям:
$$t = \frac{d_a}{d_a + d_b}, \quad
\text{pos} = (1 - t) \cdot A + t \cdot B$$

\begin{lstlisting}
let interpolate_position ~find_beacon (top2 : (string * float) list) =
  match top2 with
  | (b_a, rssi_a) :: (b_b, rssi_b) :: _ ->
    (match find_beacon b_a, find_beacon b_b with
     | Some (xa, ya, ha), Some (xb, yb, _) ->
       let d_a = distance_from_rssi rssi_a in
       let d_b = distance_from_rssi rssi_b in
       let total = d_a +. d_b in
       if total = 0. then Some (xa, ya, ha)
       else
         let t = d_a /. total in
         let x = (1. -. t) *. xa +. t *. xb in
         let y = (1. -. t) *. ya +. t *. yb in
         Some (x, y, ha)
     | _ -> None)
  | _ -> None
\end{lstlisting}

\paragraph{Точность.} Низкая. Шахтёр почти никогда не лежит точно
на отрезке между двумя биконами --- он в 2D-плоскости. Без \emph{3-го}
маяка позиция геометрически не определяется однозначно: пересечение
двух окружностей с центрами в биконах даёт две точки симметрично
относительно линии $AB$, наш алгоритм просто выбирает середину между
ними.

В реальной production-системе это нужно делать \emph{трилатерацией}
на 3+ маяках: метод наименьших квадратов на пересечении трёх (и
более) окружностей. Существенно сложнее в коде (нужен solver
нелинейной системы), но даёт реалистичную позицию.

Делаем интерполяцию намеренно: на ex07 это \emph{демо}, и
интерполяция показывает идею «координаты вычисляются из чего-то, у
нас есть что-то». Газовому пайплайну для обогащения алертов важнее
\emph{наличие хоть какой-то позиции}, чем её точность ---
``приблизительно где'' лучше для спасателя чем ``неизвестно где''.

\subsubsection{Single-beacon fallback}

Что делать если в reading'ах шахтёра \emph{один} маяк или ноль?
\texttt{interpolate\_position} вернёт \texttt{None}. Но для
газового пайплайна это плохо: первый алерт мы можем эмитить вообще
без позиции, и хотелось бы хоть какую-то указать.

\texttt{single\_beacon\_position} --- вспомогательная функция:
берёт reading с самым сильным RSSI и возвращает координаты
\emph{этого маяка}. То есть «шахтёр где-то рядом с B1, точнее
сказать не можем». Точность ниже чем у интерполяции, но это лучше
чем \texttt{None}.

\begin{lstlisting}
let single_beacon_position ~find_beacon (readings : (string * float) list) =
  match readings with
  | [] -> None
  | rs ->
    let (b, _) = List.fold_left
      (fun ((_, best) as acc) ((_, r) as cur) ->
         if r > best then cur else acc)
      (List.hd rs) (List.tl rs) in
    find_beacon b
\end{lstlisting}

Эта функция используется в газовом пайплайне (раздел~\ref{ssec:gas})
для обогащения координат \emph{до} того как первое окно RSSI
закрылось --- сразу после первого raw RSSI-пакета шахтёра.

\subsubsection{Собирая всё вместе}

Финальный пайплайн:

\begin{lstlisting}
let median_rssi ?(find_beacon = Domain.beacon_coords)
    (source : packet Mf_event.t Stream.t)
    : location Mf_event.t Stream.t =
  source
  |> Pipe.dedup (module ByLamp)
       ~rule:(fun p -> string_of_int p.ts)
       ~cooldown:(seconds 120)
  |> Pipe.flat_map (fun p -> ... (* expand to readings *))
  |> Pipe.event_time ~lateness:(seconds 1)
  |> Pipe.window_agg_keyed ~by:(fun r -> r.r_lamp)
       ~allowed_lateness:(seconds 60)
       (Pipe.sliding (seconds 60) (seconds 5))
       Agg.(group_by ~key:... ~inner:median |> map top_k_by_2)
  |> map_to_location_with_position ~find_beacon
\end{lstlisting}

Шесть операторов, каждый делает ровно одну вещь. Можно прочитать
сверху вниз и понять полный жизненный цикл события. Это и есть
обещание stream processing fluent-DSL: пайплайн читается как описание
\emph{бизнес-преобразования}, без runtime-деталей про state и
таймеры.

\paragraph{Что эмитит этот пайплайн.} Поток событий
\texttt{location Mf\_event.t}: \texttt{Data} при закрытии окна,
\texttt{Retract} плюс новый \texttt{Data} при опоздавшем пакете в
\texttt{allowed\_lateness}. \texttt{Watermark}'ы продвигаются по
event-time.

Sink через \texttt{Pipe.materialize} сворачивает поток в \emph{финальную
таблицу} «(шахтёр, конец окна) $\to$ location» --- retract'ы
автоматически отменяют старые значения. Это позволяет диспетчерскому
пульту показывать \emph{актуальную картину}, не флудя на каждом
обновлении.

Дальше --- самый интересный пайплайн, газовый. У него три потока на
входе, и retract-семантика не из коробки оператора окон, а написана
вручную.

% Раздел 4 — Retract как сквозной механизм
\section{Retract как сквозной механизм}\label{sec:retract}
\textit{Этот раздел ещё не написан.}

% Раздел 5 — Производительность
\section{Производительность}\label{sec:performance}
\textit{Этот раздел ещё не написан.}

% Раздел 6 — Выводы
\section{Выводы}\label{sec:conclusions}
\textit{Этот раздел ещё не написан.}

% ─────────────────────────────────────────────────────────────────────

\end{document}
